\section{Discussion}
\label{sec:limits}
This section interprets the results against the design contract and then states
what the evidence does not cover.

\subsection{Interpretation}
The implementation routes selected task and direct-call paths through a
shared named-operation boundary. Tests verify the applicable task, approval,
receipt, readback, and sample links across two different simulator mappings;
they do not establish route equivalence.

The blueprint replay adds protocol access, heterogeneous MES data, channel
construction, acquisition, and one approved simulator write. It was not
repeated with changed profiles or bindings, so it supports only the documented
blueprint workflow.

Route-study samples ranged from 45.69 to 46.69\,$\mu$m for cast film and
674.40 to 679.60\,$^{\circ}$C for annealing, all outside their target bands.
Each was a single probe 20\,s after an operation; pre-task values were already
out of band, so these readings establish neither a response curve nor
convergence.

The four-scenario campaign reports mission-terminal values inside the
scenario windows for injection molding, wastewater treatment, annealing, and
BOPET film. Each has one deterministic scripted trajectory; for wastewater
and annealing, the terminal mission value differs from the final structured
sample, as Fig.~5 marks. This shows scenario-specific feedback, not
repeatability or route-level process attainment.

This distinction is central to the framework: readback confirms that a
setpoint reached the configured device, while the later process sample records
whether the target was met. Keeping these as separate linked records avoids
equating command completion with product attainment. The tests validate this
reporting behavior, not improved control quality.

All route-study runs used SQLite fallback because MQTT and Timescale were
unavailable; they therefore do not validate the online DAQ pipeline. The
onboarding replay used
isolated online MQTT/Timescale services. Together, the tests cover two
configured route mappings and one separate blueprint-driven onboarding
workflow.

Under mock state, comparing the approved line--product--recipe--run tuple
blocked product/run mutations that the retained checks alone did not reject.
This is a software-dispatch result; synthetic approval does not measure human
review performance \cite{parasuramanmanzey2010complacency,parasuraman2000model,bainbridge1983ironies}.

\subsection{Scope and Limitations}
Only the reported adapters and scene configurations were exercised. The
single blueprint replay does not test repeated onboarding, profile/binding
reuse, unseen scenes, or physical transfer; the runs used working-tree builds without pinned
source commits. The study also lacks a platform baseline and does not establish
adapter interchangeability, online DAQ operation in route trials, or comparative
efficiency. Such claims require matched deployments and shared integration
effort or record-completeness measures.

The route study used three sequential tasks and one post-trial sample per task
on one simulator per scene; the four-scenario campaign used one trajectory per
scene. These are not independent replicates and cannot establish a response
curve or repeatability. Follow-up work should vary profiles, repeat trials
with pre/post measurements, and compare a defined alternative integration
path. Operation restrictions do not establish security
\cite{greshake2023injection,liu2024pinjection,iec62443-3-3} or replace
independent PLC/safety-instrumented system (SIS) safeguards \cite{iec61511}.

Scope and channel-policy refusals are not uniformly persisted; auditability
is limited to the records and denial classes described in Section~IV.

\section{Conclusion}
\label{sec:conclusion}
AgentWorkShop brings task adapters, production context, named operations,
semantic I/O, configured protocol drivers, and service-owned records into a
modular task-to-device framework. Tested task routes and direct calls used a
shared operation boundary across cast-film and annealing simulator mappings
while preserving route-specific records. A separate replay completed one
documented coating-line onboarding configuration with online DAQ checks and a
simulator-confirmed smoke write. A four-scenario simulator campaign recorded
scripted control sequences and mission-terminal values within scenario target
windows, while route-study samples remained below target. A mock-state ablation
isolated the effect of an approval-bound context check. The study establishes
selected integration, scenario-specific simulator behavior, one simulated
onboarding instance, and software-contract behavior; it does not establish
repeatability, route equivalence, process improvement, configuration reuse,
physical transfer, general portability, or comparative platform benefit.
